% Source PDF pages 104--105; manual pages 2-75--2-76.
\subsection{Initial Impact Point Calculations}\label{sec:iip-calculations}

The initial impact point (IIP) is estimated by integrating a ballistic trajectory from
the vehicle's current state to impact. Propulsive forces are removed, and the
aerodynamic model is reduced to a constant drag coefficient specified through an
input ballistic coefficient $\beta_{\mathrm{iip}}$. The IIP approximates the impact
location following a propulsion-system failure and is therefore useful for range
safety analysis.

The IIP trajectory uses the equations of motion in \cref{sec:equations-of-motion}.
Its acceleration contributions are gravity, computed as described in
\cref{sec:gravity-model}, and aerodynamic drag; propulsive acceleration is zero.
The source derives the aerodynamic acceleration as
\begin{equation}
\begin{aligned}
  \vec a_{\mathrm{aero}}
  &=\frac{\vec F_{\mathrm{aero}}}{m}
   =\left(\frac{C_DqS_{\mathrm{ref}}}{m}\right)\hat x_{V_w}\\
  &=\left(\frac{C_DS_{\mathrm{ref}}qg}{W}\right)\hat x_{V_w},
\end{aligned}
\label{eq:iip-aerodynamic-acceleration-coefficients}
\end{equation}
which becomes
\begin{equation}
  \vec a_{\mathrm{aero}}
  =\left(\frac{qg}{\beta_{\mathrm{iip}}}\right)\hat x_{V_w},
  \label{eq:iip-aerodynamic-acceleration-ballistic}
\end{equation}
and, using dynamic pressure,
\begin{equation}
  \vec a_{\mathrm{aero}}
  =\frac{g\left(0.5\rho\lVert\vec V_{W\oplus}\rVert^2\right)}
         {\beta_{\mathrm{iip}}}\hat x_{V_w}.
  \label{eq:iip-aerodynamic-acceleration-dynamic-pressure}
\end{equation}
As in \cref{sec:velocity-coordinate},
\begin{equation}
  \hat x_{V_w}
  =\frac{\vec V_{W\oplus}}{\lVert\vec V_{W\oplus}\rVert},
  \label{eq:iip-wind-velocity-unit-vector}
\end{equation}
so the acceleration is written
\begin{equation}
  \vec a_{\mathrm{aero}}
  =\frac{0.5g\rho\lVert\vec V_{W\oplus}\rVert\vec V_{W\oplus}}
         {\beta_{\mathrm{iip}}}.
  \label{eq:iip-aerodynamic-acceleration-vector}
\end{equation}

\noindent\textit{Editorial note.} Equations
\ref{eq:iip-aerodynamic-acceleration-coefficients}--
\ref{eq:iip-aerodynamic-acceleration-vector} are reproduced with the positive
$\hat x_{V_w}$ sign printed in the source. This conflicts with the drag direction in
\cref{eq:aero-force-lift-drag-side}, where drag acts along
$-\hat x_w$. The inconsistency is retained for traceability and recorded in the
editorial metadata.

TAOS integrates the IIP trajectory with Fehlberg's adaptive Runge--Kutta method.\manualref{9}
The method uses a fourth-order solution together with a fifth-order estimate to
control local error and adjust step size. This is appropriate because only the final
impact state is needed; there are no trajectory segments or print-interval output
points in the auxiliary integration.

The step size is bounded between $0.5$ and $60$ seconds while the estimated
integration error is held between $0.001$ and $0.01$ ft. Typical steps are
$10$--$20$ seconds, although the IIP calculation remains computationally
significant because it is repeated from the current trajectory state.

Impact time, longitude, and geodetic latitude are obtained directly from the
integration. IIP range and azimuth are measured from the tangent-plane origin and
computed with Sodano's inverse method in \cref{sec:ranges-distances}. The
function \texttt{iip\_calc} performs the integration, \texttt{iip\_derivs} evaluates
the IIP equations of motion, and \texttt{iip\_out} computes the output variables.
